\documentclass[11pt]{article}
\usepackage[margin=1in]{geometry}
\usepackage{amsmath,amssymb,amsthm,mathtools}
\usepackage{enumitem}
\usepackage{booktabs,longtable}
\usepackage{array}
\usepackage[hidelinks]{hyperref}
\newcommand{\cref}[1]{\ref{#1}}
\newcommand{\Cref}[1]{\ref{#1}}

\newtheorem{theorem}{Theorem}[section]
\newtheorem{lemma}[theorem]{Lemma}
\newtheorem{proposition}[theorem]{Proposition}
\newtheorem{corollary}[theorem]{Corollary}
\newtheorem{definition}[theorem]{Definition}
\newtheorem{convention}[theorem]{Convention}
\theoremstyle{remark}
\newtheorem{remark}[theorem]{Remark}

\newcommand{\C}{\mathcal C}
\newcommand{\F}{\mathcal F}
\newcommand{\G}{\mathcal G}
\newcommand{\E}{\mathbb E}
\newcommand{\Prob}{\mathbb P}
\newcommand{\HH}{\mathrm H}
\newcommand{\KL}{\mathrm D}
\newcommand{\Split}{\operatorname{Split}}
\newcommand{\Law}{\operatorname{Law}}
\newcommand{\Coll}{\operatorname{Coll}}
\newcommand{\supp}{\operatorname{supp}}
\newcommand{\rank}{\operatorname{rank}}
\newcommand{\Old}{\operatorname{Old}}
\newcommand{\Fresh}{\operatorname{Fresh}}
\newcommand{\State}{\mathsf S}
\newcommand{\Ret}{\mathsf R}
\newcommand{\Val}{\mathsf V}
\newcommand{\Eq}{\mathsf E}
\newcommand{\pr}{\mathrm{pr}}
\newcommand{\TV}{\operatorname{TV}}
\newcommand{\win}{\mathrm{win}}
\newcommand{\old}{\mathrm{old}}
\newcommand{\fr}{\mathrm{fr}}
\newcommand{\scal}{\mathrm{sc}}
\newcommand{\rk}{\mathrm{rk}}

\title{Retained-State Entropy Accounting and Fixed Windows for Three-Petal Sunflowers}
\author{Yuli Li \and Jingping Yang \and Renhao Zhang}
\date{May 2026}

\begin{document}
\maketitle

\begin{abstract}
We develop a retained-state coordinate-collision entropy method for three-petal-sunflower-free set families.  Starting from set families of size at most $k$, we reduce to balanced two-collision codes.  We prove two terminal descent principles: projected value-shadow descent and equality-shadow descent.  Equality shadows allow a public event $X_J=Y_J$ to terminate without exposing the common alphabet value.  We also prove a fixed labelled hash-window copying lemma: once a bounded public key and a fixed number of low-collision coordinates for one sample label have been obtained, three conditioned copies return a near-sunflower block at cost $O(B_{\rm win}+h)$, independent of the alphabet sizes.

The retained-state ledger is the bookkeeping mechanism.  Internal equality patterns used to analyse old-supported sparse residual witnesses are not permanent transcript variables.  Before any nonterminal child is entered, every random old coordinate block that will be used later is either terminalized or encoded as an increment of capped old-coordinate marks.  Future choices are recomputed from the conditional law on the quotient event determined by the retained state.  Thus repeated renewal of sample signatures is charged through a final finite state of entropy $O(|U|)$, rather than through the sum of all internal possible-support pattern exposures.

The proof proceeds through a retained-state assembly theorem.  The local transition theorem in \Cref{sec:local-transition} is organized as a finite list of mode-by-mode alternatives: every close-pair, product-KL, hash, residual, old-window, reservoir, and scale-birth transition is represented as terminal data, monotone retained data, charged pruning, or analysis data erased before continuation.  Combining these local kernels with the retained entropy ledger gives a constant-base recurrence for balanced two-collision codes, and hence for three-petal-sunflower-free set families.
\end{abstract}

\section{Statement of the reduction}

A three-petal sunflower is a triple of distinct sets $A,B,C$ such that
\[
        A\cap B=A\cap C=B\cap C.
\]
A family has size at most $k$ if each member has cardinality at most $k$.

\begin{definition}[Balanced two-collision code]
A code
\[
        \C\subseteq \prod_{i=1}^{k}\Omega_i
\]
is a balanced two-collision code of length $k$ if for every three distinct codewords $x,y,z\in\C$ there is a coordinate $i$ for which exactly two of $x_i,y_i,z_i$ are equal.
\end{definition}

\begin{lemma}[Padding]
Let $\F$ be a family of sets of size at most $k$ with no three-petal sunflower.  Then there is a $k$-uniform family $\F^+$ with $|\F^+|=|\F|$ and no three-petal sunflower.
\end{lemma}

\begin{proof}
For each $A\in\F$ add $k-|A|$ private dummy elements, disjoint from the original ground set and from the dummy elements used for all other members.  If $A^+,B^+,C^+$ formed a sunflower, no private dummy element could lie in an intersection of two distinct padded sets.  Hence the pairwise intersections of $A^+,B^+,C^+$ equal those of $A,B,C$, and $A,B,C$ would be a sunflower in $\F$.
\end{proof}

\begin{lemma}[Random colouring]
Let $\F$ be a $k$-uniform family with no three-petal sunflower.  Then there is a balanced two-collision code $\C$ of length $k$ such that
\[
        |\C|\ge \frac{k!}{k^k}|\F|.
\]
\end{lemma}

\begin{proof}
Colour the ground set independently and uniformly with colours $1,\ldots,k$.  A fixed $k$-set is colourful with probability $k!/k^k$, so some colouring leaves a colourful subfamily $\F_0$ of size at least $(k!/k^k)|\F|$.  Encode each $A\in\F_0$ by the ordered list of its unique element of each colour.  If three encoded words had no coordinate with exactly two equal symbols, then each colour class would contribute either one element common to all three sets or three different elements.  The three pairwise intersections of the sets would then be equal, contradicting sunflower-freeness.  Thus the encoded family is a balanced two-collision code.
\end{proof}

\begin{corollary}[Code reduction]\label{cor:code-reduction}
If every balanced two-collision code of length $k$ has size at most $D^k$, then every family of sets of size at most $k$ with no three-petal sunflower has size at most $(eD)^k$.
\end{corollary}

\begin{proof}
Combine the previous two lemmas and use $k!\ge (k/e)^k$.
\end{proof}

\section{Terminal shadows}

All random variables are finite-valued, and all entropies are natural logarithms.  Let $X^1,\ldots,X^r$ be independent samples from the uniform law on a balanced two-collision code.

\begin{definition}[Projected value shadow]
A projected value shadow is a finite list
\[
        \theta=((J_s,a_s))_{s\in S},
\]
where $\emptyset\ne S\subseteq [r]$, $\emptyset\ne J_s\subseteq [k]$, and $a_s\in\prod_{j\in J_s}\Omega_j$.  An event $E$ is compatible with $\theta$ if
\[
        E\subseteq \bigcap_{s\in S}\{X^s_{J_s}=a_s\}.
\]
Define
\[
        T(\theta)=\sum_{s\in S}|J_s|,\qquad R(\theta)=|S|.
\]
\end{definition}

\begin{theorem}[Projected value-shadow descent]\label{thm:value-shadow}
Assume inductively that every balanced two-collision code of length below $k$ has size at most $D^j$ in length $j$.  Let $\C$ be a length-$k$ balanced two-collision code, and suppose an event $E$ compatible with a nonempty projected value shadow $\theta$ satisfies
\[
        \Prob(E)\ge \exp(-\lambda T(\theta)-cR(\theta)).
\]
If $\log D\ge \lambda+c$, then $|\C|\le D^k$.
\end{theorem}

\begin{proof}
Let $M=|\C|$ and let $\mu$ be uniform on $\C$.  For each $s\in S$, set
\[
        \eta_s=\mu\{x:x_{J_s}=a_s\}.
\]
Independence gives $\Prob(E)\le\prod_s\eta_s$.  For each $s$, the cylinder $\{x:x_{J_s}=a_s\}$, after deleting the coordinates $J_s$, remains a balanced two-collision code of length at most $k-|J_s|$.  Hence $\eta_s M\le D^{k-|J_s|}$.  Multiplying over $s$ gives
\[
        M^{R(\theta)}\Prob(E)\le D^{R(\theta)k-T(\theta)}.
\]
Using the lower bound for $\Prob(E)$,
\[
        M^{R(\theta)}\le \exp(\lambda T(\theta)+cR(\theta))D^{R(\theta)k-T(\theta)}.
\]
Since $T(\theta)\ge R(\theta)$ and $\log D\ge \lambda+c$, the exponential factor is at most $D^{T(\theta)}$.  Therefore $M^{R(\theta)}\le D^{R(\theta)k}$.
\end{proof}

\begin{definition}[Equality shadow]
An equality shadow is a public nonempty coordinate block $J$ and two sample labels $s\ne t$, recording the event
\[
        X^s_J=X^t_J.
\]
It records no alphabet value.
\end{definition}

\begin{theorem}[Equality-shadow descent]\label{thm:equality-shadow}
Assume the same induction hypothesis as in \Cref{thm:value-shadow}.  Let $X,Y$ be independent uniform samples from a length-$k$ balanced two-collision code $\C$, and let $J\ne\emptyset$ be public.  If an event $E\subseteq\{X_J=Y_J\}$ satisfies
\[
        \Prob(E)\ge \exp(-\lambda |J|-c)
\]
and $\log D\ge \lambda+c$, then $|\C|\le D^k$.
\end{theorem}

\begin{proof}
Let $M=|\C|$.  For each value vector $a\in\prod_{j\in J}\Omega_j$, put $\C_a=\{x\in\C:x_J=a\}$ and $\eta_a=|\C_a|/M$.  Since $E\subseteq\{X_J=Y_J\}$,
\[
        \Prob(E)\le \sum_a \eta_a^2.
\]
Each nonempty $\C_a$, after deleting $J$, is a balanced two-collision code of length at most $k-|J|$, so $|\C_a|\le D^{k-|J|}$.  Thus
\[
        \sum_a\eta_a^2=\frac{1}{M^2}\sum_a|\C_a|^2
        \le \frac{D^{k-|J|}}{M}.
\]
The lower bound on $\Prob(E)$ gives $M\le \exp(\lambda |J|+c)D^{k-|J|}\le D^k$.
\end{proof}

\section{Retained-state proof trees}

The proof tree will use local analysis partitions.  The essential point is that analysis partitions need not become permanent transcript variables.

\begin{definition}[Retained transition]
At a nonterminal state $\State$, a local transition consists of:
\begin{enumerate}[label=\textup{(\roman*)}]
\item a finite or countable analysis partition $\mathcal A$ of the current state event;
\item for each analysis atom, either a terminal output or a retained label $Y$;
\item a quotient rule which merges all analysis atoms having the same retained label $Y$.
\end{enumerate}
The child law for the label $Y=y$ is the conditional law on the union of all analysis atoms with retained label $y$.  Future selectors are required to be deterministic functions of this quotient law, the fixed thresholds, and the retained label.  They may not depend on the discarded analysis atom.
\end{definition}

\begin{definition}[Admissible retained data]\label{def:retained-data}
The retained data of a nonterminal state consists only of:
\begin{enumerate}[label=\textup{(\roman*)}]
\item the current conditional law on the retained active sample labels;
\item public coordinate blocks and active scale labels;
\item the old support $U$;
\item capped old-coordinate marks $M:U\to\{0,\ldots,Q\}$;
\item bounded old queues and active blocks, each encoded by finitely many per-coordinate marks on $U$;
\item pending hash tokens, at most one on the unique residual state immediately returned by a fixed hash window;
\item terminal value shadows, equality shadows, and pruning losses already recorded.
\end{enumerate}
No discarded equality pattern, witness-label choice, or old block order belongs to the retained state unless it is encoded by one of the finite objects above.
\end{definition}

\begin{lemma}[Retained-state entropy]\label{lem:retained-entropy}
Let a finite retained-state tree terminate with terminal object
\[
        Z_*=(\Theta_{\rm val},\Theta_{\rm eq},S_*)
\]
where $\Theta_{\rm val}$ is the value-shadow data, $\Theta_{\rm eq}$ is equality-shadow data, and $S_*$ is the final retained non-value state.  If every future selector after a quotient is measurable with respect to the quotient law and the retained label, then
\[
        \HH(\Theta_{\rm val},\Theta_{\rm eq})\le \HH(Z_*).
\]
Moreover, if the retained output at a transition has conditional entropy at most $c_v$ plus pruning loss $\ell^{\pr}_v$, then
\[
        \HH(Z_*)+\E L_{\pr}\le \E\sum_{v\ {\rm reached}} c_v+\E L_{\pr}.
\]
Analysis-only variables which are quotiented before the next state do not appear in the right-hand side.
\end{lemma}

\begin{proof}
The terminal shadows are deterministic functions of $Z_*$.  For the second assertion, reveal only the retained labels along the reached path, not the full analysis atoms.  The chain rule gives the sum of the conditional entropies of the retained labels.  The quotient measurability condition ensures that later transition kernels are functions of the already revealed retained labels and their quotient laws; hence discarded atoms are not needed to reconstruct the distribution of future retained labels.  Adding pruning losses is an ordinary normalization account and does not require retaining the discarded branch.
\end{proof}

\section{Elementary local tools}

\subsection{Collision and DRC}

For a one-sample law $\nu$ on a code and a coordinate $i$, let
\[
        \Coll_i(\nu)=\Prob_{X,Y\sim\nu}[X_i=Y_i].
\]
For three samples $X,Y,Z$, let $\Split_J(X,Y,Z)$ be the set of coordinates in $J$ at which exactly two of $X_i,Y_i,Z_i$ are equal.

\begin{lemma}[Collision controls split]\label{lem:collision-split}
For independent $X,Y,Z\sim\nu$,
\[
        \E |\Split_J(X,Y,Z)|\le 3\sum_{i\in J}\Coll_i(\nu).
\]
Consequently, if the average collision on $J$ is at most $\beta$, then
\[
        \Prob\bigl[|\Split_J(X,Y,Z)|\le 6\beta |J|\bigr]\ge \frac{1}{2}.
\]
\end{lemma}

\begin{proof}
For one coordinate, the probability that at least two of $X_i,Y_i,Z_i$ are equal is at most $3\Prob[X_i=Y_i]=3\Coll_i(\nu)$.  Sum over $i$ and apply Markov's inequality.
\end{proof}

\begin{lemma}[A DRC row block]\label{lem:drc}
Let $(\Omega,\Prob)$ be a probability space and let $E_i$, $i\in I$, be events with average mass at least $\delta$.  For every $1\le s\le \delta |I|/4$ there is a block $J\subseteq I$, $|J|=s$, such that
\[
        \Prob\Bigl[\bigcap_{j\in J}E_j\Bigr]\ge \frac{\delta}{2}\left(\frac{\delta}{4}\right)^s.
\]
If a fixed public order on $s$-subsets of $I$ is given, the first such $J$ is public.
\end{lemma}

\begin{proof}
Choose $s$ indices independently from $I$ according to the distribution proportional to $\Prob(E_i)$, then use the standard dependent-random-choice averaging argument; equivalently,
\[
        \E_{\omega}\binom{d(\omega)}{s}\ge \binom{\delta |I|/2}{s}\Prob[d(\omega)\ge \delta |I|/2]
\]
after discarding the set on which the degree $d(\omega)$ is below $\delta |I|/2$.  The displayed bound follows from the usual estimate for $s\le \delta |I|/4$.  The public choice is by the fixed tie-break.
\end{proof}

\subsection{Positive KL scans}

\begin{lemma}[Positive tail]\label{lem:positive-tail}
If $p,q$ are laws on a finite set and $D(p\|q)>0$, then the set $S_+=\{a:p(a)>q(a)\}$ has positive $p$-mass.  If $p(S_+)<\rho\le1/2$, discarding $S_+$ costs pruning loss at most $2\rho$ on the retained complement.
\end{lemma}

\begin{proof}
If $p(a)\le q(a)$ for all $a$, then $p=q$ and the divergence is zero.  The pruning bound is $-\log(1-x)\le 2x$ for $0\le x\le1/2$.
\end{proof}

\begin{lemma}[Heavy-or-light one-coordinate routing]\label{lem:heavy-light}
Let $p$ be a finite law and fix $0<\rho,\tau\le1$.  Put
\[
        H=\{a:p(a)\ge \rho\tau\},\qquad L=H^c.
\]
Then the branches $a\in H$ are terminal value pins of conditional mass at least $\rho\tau$ and there are at most $(\rho\tau)^{-1}$ of them.  If $p(L)<\rho$, discarding $L$ costs pruning loss at most $2\rho$.  If $p(L)\ge\rho$, then
\[
        \Coll(p(\cdot\mid L))\le \tau.
\]
\end{lemma}

\begin{proof}
The number of heavy atoms is at most $(\rho\tau)^{-1}$.  If $p(L)\ge\rho$, then
\[
        \Coll(p(\cdot\mid L))=\sum_{a\in L}\left(\frac{p(a)}{p(L)}\right)^2
        \le \frac{\max_{a\in L}p(a)}{p(L)}\le \tau.
\]
\end{proof}


\begin{lemma}[One-coordinate collision exits]\label{lem:collision-exit}
Let $p,q$ be one-coordinate laws and let $0<\rho\le 1/2$.  If
\[
        \sum_a p(a)q(a)>\gamma_0,
\]
then either $\max_a p(a)>\gamma_0^2$ or $\max_a q(a)>\gamma_0^2$.  If
\[
        \Coll(p)=\sum_a p(a)^2>\alpha,
\]
then $\max_a p(a)>\alpha$.  Consequently, whenever $\rho\tau\le\min(\alpha,\gamma_0^2)$, each cross-collision or high-side-collision branch is routed through one actual sample-coordinate law $r$: heavy atoms of $r$ terminalize as value pins, a light complement of mass below $\rho$ is charged pruning, and a light complement of mass at least $\rho$ gives a low-collision actual hash coordinate.
\end{lemma}

\begin{proof}
By Cauchy's inequality,
\[
        \sum_a p(a)q(a)\le \sqrt{\Coll(p)\Coll(q)}.
\]
If both maxima were at most $\gamma_0^2$, then both collisions would be at most $\gamma_0^2$, contradicting the strict cross-collision inequality.  The high-side statement follows from $\Coll(p)\le\max_a p(a)$.  Choose $r=p$ or $r=q$ by the fixed public tie-break among the laws having a sufficiently large atom, and then apply \Cref{lem:heavy-light}.  The selected law is an actual coordinate law of a named sample label, so its nonterminal low-collision branch is eligible for a labelled hash window.
\end{proof}

\begin{proposition}[Actual-coordinate positive-KL scan]\label{prop:positive-kl-scan}
Let $V=(V_1,\ldots,V_m)$ be a public list of actual sample-coordinate variables.  Let $P\ll Q$ be laws on the product space of $V$ with $D(P\|Q)>0$.  Fix $\lambda>0$, $0<\rho\le 1/2$, and $0<\tau<1$.  The following retained transition is exhaustive.  Along a realised branch scan the public list.  If a positive one-coordinate KL stop occurs before any prefix whose $P$-mass is below $e^{-\lambda i}$, apply \Cref{lem:positive-tail} and \Cref{lem:heavy-light} to that actual coordinate under the paid prefix.  If a prefix first crosses the threshold, apply \Cref{lem:heavy-light} to the crossing coordinate under the previous amortized prefix.  If neither occurs, terminalize the full prefix.

Every output is one of:
\begin{enumerate}[label=\textup{(\roman*)}]
\item an amortized terminal value shadow prefix with cost at most $\lambda$ times its length;
\item a one-coordinate terminal value pin of conditional mass at least $\rho\tau$;
\item a charged pruning branch not inherited by a nonterminal child;
\item a labelled low-collision hash coordinate for an actual sample label, under a paid set event of cost at most $\log(1/\rho)$.
\end{enumerate}
No reference-law coordinate, pair alphabet, or erased analysis atom is retained.
\end{proposition}

\begin{proof}
Use the chain rule for relative entropy.  The conditional one-coordinate divergences are public after the realised prefix has been paid.  If a positive divergence is found, \Cref{lem:positive-tail} gives a positive tail.  If the tail has large mass, \Cref{lem:heavy-light} gives either a terminal value atom or a low-collision conditional law.  If the tail has small mass, prune it and apply the same heavy/light routing to the retained complement.  The prefix crossing case is exactly the same but the previous prefix is amortized by definition of first crossing.  If no stop occurs, the full prefix has pointwise mass at least $e^{-\lambda m}$ and is terminalized.  Each nonterminal low-collision output is a conditional law of one of the actual coordinates $V_t$; the reference law $Q$ is not passed to the child state.
\end{proof}

\section{Fixed labelled hash windows}

For a key $K$ with law $p$, write
\[
        H_3(K)=-\frac{1}{2}\log\sum_t p(t)^3.
\]

\begin{lemma}[Random-key copying]\label{lem:random-copy}
Let $K$ be a finite key.  Conditional on $K=t$, let $\nu_t$ be a one-sample law.  Suppose that for every $t$ there is an event $F_t$ for three independent draws from $\nu_t$ with $\nu_t^3(F_t)\ge\rho$.  Then three independent copies of the stopped experiment have probability at least
\[
        \rho\sum_t\Prob[K=t]^3=\rho\exp(-2H_3(K))
\]
of having the same key and satisfying the corresponding event.
\end{lemma}

\begin{proof}
The desired mass is $\sum_t\Prob[K=t]^3\nu_t^3(F_t)$.
\end{proof}

\begin{theorem}[Fixed labelled low-collision window]\label{thm:fixed-window}
Fix $h\ge1$ and $0<\beta<1/12$.  At a retained state suppose one public sample label has a stopped hash key $K$ with $\HH(K)\le B_{\rm win}$, and, for each key value $t$, a public set $H(t)$ of at least $h$ surviving coordinates such that under the conditional one-sample law $\nu_t$:
\begin{enumerate}[label=\textup{(\roman*)}]
\item $\Coll_i(\nu_t)\le\beta$ for all $i\in H(t)$;
\item every codeword atom has $\nu_t$-mass below $\beta$.
\end{enumerate}
Then three conditioned copies produce, with probability at least $\exp(-O(B_{\rm win}+h))$, a distinct triple whose split set inside $H(t)$ has size at most $6\beta h$.  The exact split exposure and the copy tax are $O(B_{\rm win}+h)$, with no dependence on alphabet size or on $|\C|$.
\end{theorem}

\begin{proof}
For fixed $t$, \Cref{lem:collision-split} gives split size at most $6\beta h$ with probability at least $1/2$.  The probability that three independent $\nu_t$-samples are not all distinct is at most $3\sum_x\nu_t(x)^2\le3\beta$.  Hence the good event has probability at least $1/2-3\beta>0$.  Apply \Cref{lem:random-copy}.  Conditioning on any good-key event of mass bounded below by a constant changes $\HH(K)$ and $H_3(K)$ by only a constant factor.  The exact split exposure has at most $2^h$ possibilities, or fewer under a sparse threshold.  Therefore the total cost is $O(B_{\rm win}+h)$.
\end{proof}

\begin{convention}[Hash-token rule]\label{conv:hash-token}
A nonterminal fixed-window return creates at most one pending token of value $C_{\rm hash}=O(B_{\rm win}+h)$.  The returned residual block must be processed until the first terminal, fresh, old, scale, or birth hard exit.  The token is assigned to that exit before any child state or descendant hash window is entered.  Bounded-support hash exits are classified immediately as fresh or old and create no pending token.
\end{convention}


\section{Close-pair and product-KL retained kernels}\label{sec:close-pair}

A close-pair state carries a public block $J$ and active samples with, after relabelling,
\[
        X_J=Y_J=A_J,
\tag{*}
\]
and a third sample $Z$ whose values on $J$ are denoted $B_J$, with $A_j\ne B_j$ on the branch.

\begin{lemma}[Conditioned product factorization]\label{lem:conditioned-product}
For each $j\in J$, let $p_j,q_j$ be probability laws on $\Omega_j$, and put
\[
        \gamma_j=\sum_a p_j(a)q_j(a)<1.
\]
Let $Q=\bigotimes_{j\in J}(p_j\otimes q_j)$ on pairs $(A_J,B_J)$ and let
\[
        E=\{A_j\ne B_j\text{ for every }j\in J\}.
\]
Then $Q(\cdot\mid E)$ factorizes as
\[
        Q_E=\bigotimes_{j\in J} Q_j^*,
        \qquad
        Q_j^*(a,b)= \frac{p_j(a)q_j(b)\mathbf 1_{a\ne b}}{1-\gamma_j}.
\]
\end{lemma}

\begin{proof}
Under $Q$ the coordinate-pair events $\{A_j\ne B_j\}$ are independent.  Thus $Q(E)=\prod_j(1-\gamma_j)$, and division by this product gives the displayed factorization.
\end{proof}

\begin{lemma}[Product transfer to near-sunflowers]\label{lem:product-near-sunflower}
Let $P=\Law(A_J,B_J)$ be supported on $A_j\ne B_j$ for every $j\in J$.  Let $p_j=\Law(A_j)$, $q_j=\Law(B_j)$, and let $Q_E$ be the conditioned product law of \Cref{lem:conditioned-product}.  Assume
\[
        \KL(P\|Q_E)\le \epsilon,
        \qquad
        \gamma_j\le\gamma_0<1\quad(j\in J),
\]
and
\[
        \frac{1}{|J|}\sum_{j\in J}\Coll(p_j)\le\alpha.
\]
If $A^{(1)},A^{(2)},A^{(3)}$ are independent samples from the true $A$-marginal $P_A$, then
\[
        \Prob\left[|\Split_J(A^{(1)},A^{(2)},A^{(3)})|
        \le \frac{6\alpha |J|}{(1-\gamma_0)^2}\right]
        \ge \frac{1}{2}-3\sqrt{\epsilon/2}.
\]
\end{lemma}

\begin{proof}
Under $Q_E$, the $A_j$-marginal is
\[
        p_j^*(a)=\frac{p_j(a)(1-q_j(a))}{1-\gamma_j}.
\]
Since $1-q_j(a)\le1$ and $1-\gamma_j\ge1-\gamma_0$,
\[
        \Coll(p_j^*)=\sum_a \frac{p_j(a)^2(1-q_j(a))^2}{(1-\gamma_j)^2}
        \le \frac{\Coll(p_j)}{(1-\gamma_0)^2}.
\]
Thus three independent samples from the $A$-marginal of $Q_E$ have expected split count at most $3\alpha |J|/(1-\gamma_0)^2$ by \Cref{lem:collision-split}; Markov gives probability at least $1/2$ for the displayed split bound.  Pinsker's inequality gives $\TV(P,Q_E)\le\sqrt{\epsilon/2}$.  Total variation cannot increase under marginalization, and the distance between threefold product laws is at most three times the one-sample distance.  Hence the same event has probability at least $1/2-3\sqrt{\epsilon/2}$ under $P_A^3$.
\end{proof}

\begin{lemma}[Projection near-sunflowers lift to code triples]\label{lem:projection-lift}
Let $\mu$ be a law on a code $\C$, let $\pi_J$ be projection to $J$, and let $\nu=(\pi_J)_*\mu$.  If a projected event $F\subseteq(\pi_J\C)^3$ has $\nu^3(F)\ge\rho$, then at least one of the following holds:
\begin{enumerate}[label=\textup{(\roman*)}]
\item the first codeword $x$ in the fixed public order satisfying $\mu(x)\ge\rho/6$ is retained as a terminal codeword atom;
\item for independent $X,Y,Z\sim\mu$, the event $(\pi_JX,\pi_JY,\pi_JZ)\in F$ occurs with $X,Y,Z$ distinct with probability at least $\rho/2$.
\end{enumerate}
\end{lemma}

\begin{proof}
The projections of independent samples from $\mu$ are independent samples from $\nu$, so the projected event has probability exactly $\rho$ before imposing distinctness.  Let $\Delta=\Prob[X=Y]=\sum_x\mu(x)^2$.  The probability that three samples are not all distinct is at most $3\Delta$.  If $\Delta\ge\rho/6$, then $\max_x\mu(x)\ge\Delta\ge\rho/6$, giving the terminal atom.  Otherwise $3\Delta<\rho/2$, and the distinct lift has mass at least $\rho/2$.
\end{proof}

\begin{proposition}[Close-pair master alternatives]\label{prop:close-pair}
Let $P=\Law(A_J,B_J)$, let $p_j=\Law(A_j)$, $q_j=\Law(B_j)$, and $\gamma_j=\sum_a p_j(a)q_j(a)$.  Fix constants $h,\alpha,\epsilon,\gamma_0$ with $\gamma_0<1$ and
\[
        \rho_*=\frac{1}{2}-3\sqrt{\epsilon/2}>0.
\]
At least one of the following retained exits holds:
\begin{enumerate}[label=\textup{(\arabic*)}]
\item $\HH(A_J)\le h|J|$, in which case an entropy atom of $A_J$ terminalizes as a value shadow, and if the retained block is already equality-public it may terminalize as an equality shadow;
\item $\gamma_j>\gamma_0$ for some $j$, in which case \Cref{lem:collision-exit} gives a terminal pin, charged pruning, or a labelled low-collision hash coordinate;
\item all $\gamma_j<1$ and the product-KL alternative holds:
\[
        \KL(P\|Q_E)>\epsilon;
\]
this branch is routed by \Cref{prop:product-kl};
\item the average $A$-side collision on $J$ exceeds $\alpha$, giving a public DRC close-pair subblock;
\item a terminal codeword atom of mass at least $\rho_*/6$ appears;
\item a law-level near-sunflower block on $J$ returns to residual mode with exact split exposure and mass at least $\rho_*/2$.
\end{enumerate}
All nonterminal children are functions of the current quotient law and retained labels only.
\end{proposition}

\begin{proof}
If some $\gamma_j=1$, then alternative (2) holds because $\gamma_0<1$.  Thus, outside alternative (2), $Q_E$ is defined.  If alternative (1), (3), or (4) holds, use the stated routing.  In alternative (4), apply \Cref{lem:drc} to the events that two independent $A$-samples agree at coordinate $j$; the block is chosen by the fixed public order and the role information is retained or terminalized by the close-pair rules.

Assume the first four alternatives fail.  Then
\[
        \KL(P\|Q_E)\le\epsilon,
        \qquad
        \gamma_j\le\gamma_0,
        \qquad
        \frac{1}{|J|}\sum_j\Coll(p_j)\le\alpha.
\]
By \Cref{lem:product-near-sunflower}, three independent samples from the true $A$-marginal have a projected near-sunflower event of mass at least $\rho_*$.  This $A$-marginal is exactly the $J$-projection of the current one-sample law $\mu$.  Applying \Cref{lem:projection-lift} gives either alternative (5) or alternative (6).  In alternative (6) the exact split set on $J$ is exposed at the transition and the child law is the quotient law determined by that retained split exposure.
\end{proof}

\begin{proposition}[Product-KL retained kernel]\label{prop:product-kl}
In the product-KL alternative of a close-pair state, list the public block as $J=(j_1,\ldots,j_m)$ and scan the actual tuple-coordinate list
\[
        V=(A_{j_1},B_{j_1},A_{j_2},B_{j_2},\ldots,A_{j_m},B_{j_m}).
\]
Then the retained outputs are exactly those of \Cref{prop:positive-kl-scan}.  In particular, a nonterminal hash coordinate is always one actual sample-coordinate value $X^s_{j_r}$ under a paid prefix/set event.  The pair alphabet and the reference product law are not retained.
\end{proposition}

\begin{proof}
The interleaving is a bijection of the pair data $(A_J,B_J)$, so relative entropy is invariant under it.  Apply \Cref{prop:positive-kl-scan} to the pushforwards of $P$ and $Q_E$.  The selected low-collision variable is either $A_{j_r}=X_{j_r}=Y_{j_r}$ or $B_{j_r}=Z_{j_r}$, hence is an actual code coordinate for a named sample label.  Terminal prefixes and pins become multi-sample value shadows; small complements are pruning losses; large complements are paid set-events inside the corresponding labelled hash window.  The pair alphabet and $Q_E$ are used only to locate the stopping time and are not retained by the child.
\end{proof}

\section{Residual states and sparse old renewal}\label{sec:residual}

A residual state carries a public near-sunflower block and an exact split exposure.  Let $U$ be the old support.  Coordinates outside $U$ are fresh.

\begin{lemma}[Fresh support promotion]\label{lem:fresh-promotion}
At a residual state, if the deterministic fresh positive support $F$ has size greater than a fixed constant $q$, promote $F$ to $U$.  If $|F|\le q$, expose only the five-state pattern on $F$; each atom either promotes a nonempty split part and gives a fresh close-pair child, or leaves all residual split witnesses in $U$.  The cost is at most $C_{\rm fr}(q)|F|$ and is paid by the fresh potential.
\end{lemma}

\begin{proof}
The set $F$ is determined by the current quotient law.  For $|F|>q$ the promotion is public.  For $|F|\le q$ there are at most $5^q$ equality patterns; on any atom with a nonempty split part, promote that part and choose the majority split role.  If there is no fresh split coordinate on the atom, the witness is old-supported.  The transcript cost is a fixed multiple of $|F|$.
\end{proof}

\begin{proposition}[Dense residual DRC kernel]\label{prop:dense-residual}
Let $W_L(X,Y,Z)$ be the residual split witness on a public residual domain $L$.  If
\[
        \E |W_L|\ge \beta |L|,
\]
and $|L|\ge 4/\beta$, then for $s=\lfloor\beta |L|/8\rfloor$ there is a public block $J\subseteq L$, $|J|=s$, such that all coordinates of $J$ split on a subevent of mass $\exp(-O(s))$.  After exposing a three-role vector, a subblock $J'\subseteq J$ of size at least $s/3$ is a close-pair block.  If $J'$ is fresh, it is promoted; if it is old and nonterminal, its coordinate identity is retained only through old mark increments or a bounded queue.
\end{proposition}

\begin{proof}
Apply \Cref{lem:drc} to the events $\{t\in W_L\}$.  On the event that all coordinates of $J$ split, one of the three roles $XY|Z$, $XZ|Y$, $YZ|X$ occurs on at least $s/3$ coordinates.  Role exposure costs $O(s)$.  The retained output is either terminal equality/value information, a fresh promotion, or an old block whose coordinate identity is encoded as required by the old retained-state rule.
\end{proof}

\begin{proposition}[Bounded residual kernel]\label{prop:bounded-residual}
If $\E|W_L|\ge \beta |L|$ but $|L|<4/\beta$, the exact subset $W_L\subseteq L$ and its role pattern have bounded entropy depending only on $\beta$.  The branch is terminal, bounded-leaf, fresh, or old; in the old nonterminal case the actual used coordinates are recorded as old mark increments before any child is entered.
\end{proposition}

\begin{proof}
The number of possible subsets and role patterns is bounded by a constant depending only on $\beta$.  Thus the selector cost is a fixed constant.  The retained-state rule prevents the exposed subset from being carried forward unless it is encoded in the finite state.
\end{proof}

\begin{definition}[Old mark vector]
For an old support $U$, the old mark vector is a function
\[
        M:U\to\{0,1,\ldots,Q\}.
\]
Whenever an actual old coordinate block $K\subseteq U$ is used to create a nonterminal old continuation, the transition increments $M(i)$ for all $i\in K$, unless doing so would exceed $Q$, in which case the branch takes the high-revisit exit.  Active old blocks and bounded old queues may be retained, but every coordinate in them must be contained in the support of a mark increment at or before the time the block is first carried forward.
\end{definition}

\begin{lemma}[Coordinate-mark compression]\label{lem:mark-compression}
Fix $U$ and $Q$.  Along any branch satisfying the old mark rule, the entropy of the retained non-value old coordinate data, conditional on $U$, is at most
\[
        |U|\log(Q+1)+C_{\rm blk}|U|+O(1),
\]
where $C_{\rm blk}$ depends only on the bounded number of active old blocks, queue labels, role labels, and scale labels.
\end{lemma}

\begin{proof}
The vector $M$ has at most $(Q+1)^{|U|}$ possibilities.  Each retained active old block or queue is encoded by a bounded number of per-coordinate states: absent, present with one of finitely many role/scale labels, or present in a bounded queue position.  Since the number of active blocks, queue positions, roles, scales, and sample-label signatures retained at a nonterminal state is bounded by fixed constants, the number of such annotations is at most $\exp(C_{\rm blk}|U|+O(1))$.
\end{proof}

\begin{proposition}[Sparse old-supported retained kernel]\label{prop:sparse-old}
Suppose a residual branch is old-supported:
\[
        W_L(X,Y,Z)\subseteq U,
        \qquad W_L\ne\emptyset \text{ a.s.},
\]
and $\E|W_L|<\beta |L|$.  Let
\[
        O=\{u\in U:\Prob[u\in W_L]>0\},
\]
which is determined by the current quotient law.  The transition may use the five-state equality pattern on $O$ as an analysis partition.  On each positive atom the actual split set $K=W_L$ is determined; after choosing a majority role, let $K_0\subseteq K$ be the corresponding close-pair block.

Before any nonterminal child is entered, exactly one of the following happens:
\begin{enumerate}[label=\textup{(\roman*)}]
\item a terminal value or equality shadow is output;
\item registering $K_0$ would exceed the cap $Q$, and the high-revisit hard exit is taken on an actual coordinate of $K_0$;
\item the coordinates of $K_0$ are encoded as an increment of $M$, possibly also as a bounded queue or active old block; all other pattern data are quotiented.
\end{enumerate}
Thus repeated sparse renewal retains only the compressed old state of \Cref{lem:mark-compression}; no term of the form $\sum_\ell |O_\ell|$ enters the retained entropy.
\end{proposition}

\begin{proof}
The pattern on $O$ is an analysis variable, not a child selector.  It determines $K$ on each atom.  If the branch terminalizes, the only retained data are the terminal shadow and pruning loss.  If the cap is crossed, the crossing coordinate belongs to the current actual split block $K_0$, so the high-revisit exit is based on an actual close-pair revisit, not on possible support.  Otherwise increment $M$ on $K_0$ and retain only the finite role, queue, and scale labels needed for the next old-window transition.  The child law is the conditional law on the union of all atoms having the same updated retained state.  Future choices are recomputed from this law; the discarded equality pattern cannot later select a shadow or block.  The entropy bound is \Cref{lem:mark-compression}.
\end{proof}

\section{Old windows, reservoirs, and scale birth}\label{sec:old-window}


\begin{lemma}[High-revisit one-coordinate routing]\label{lem:high-revisit}
At an old close-pair occurrence suppose a coordinate $i$ has role $X_i=Y_i\ne Z_i$, and put $A_i=X_i=Y_i$ under the current conditional law.  For every $0<\rho\le1/2$ and $0<\tau<1$, the following public routing is exhaustive.  Let
\[
        H=\{a:\Prob[A_i=a]\ge\rho\tau\},\qquad L=\Omega_i\setminus H.
\]
The branches $A_i=a\in H$ are terminal one-coordinate value pins of conditional mass at least $\rho\tau$.  If $\Prob[A_i\in L]<\rho$, the light complement is a charged pruning branch of loss at most $2\rho$.  If $\Prob[A_i\in L]\ge\rho$, then after conditioning on the public event $A_i\in L$, the conditional law of the actual sample coordinate has collision at most $\tau$ and enters the labelled hash window for that sample label, provided the paid set event fits the stopped budget; otherwise the branch terminalizes on the already paid prefix.
\end{lemma}

\begin{proof}
This is \Cref{lem:heavy-light} applied to the actual one-coordinate law of $A_i$.  The coordinate identity $i$ belongs to the current actual old block which caused the cap crossing, so the chosen sample label and coordinate are retained data.  Heavy atoms terminalize, small complements are pruning, and large complements give the low-collision hash law.  No common old value vector is retained by a nonterminal child.
\end{proof}

\begin{proposition}[Fixed-scale old-window retained kernel]\label{prop:old-window}
Let $O_1,\ldots,O_R\subseteq U$ be public actual split old close-pair blocks of one dyadic scale and one public sample-label signature.  Deterministic possible-support sets used only to expose equality patterns are not members of this list.  The old-window transition has the following retained exits:
\begin{enumerate}[label=\textup{(\roman*)}]
\item a terminal value or equality shadow;
\item a high-revisit hard exit on a coordinate whose mark would exceed $Q$;
\item a low-revisit continuation in which all actual old incidences are encoded as increments of $M$ and any bounded short window is put into the bounded queue;
\item a signature-changing hard exit which first charges or registers the bounded queue and then clears it.
\end{enumerate}
No high-symbol old value vector is retained in a nonterminal child.
\end{proposition}


\begin{proof}
The persistent sample-label signature and dyadic scale are part of the public state, so every block in the list refers to the same ordered labels and the same scale.  Process the queued blocks in their public order.  If the branch elects to stop on a public close-pair block, it records either a value shadow or the equality shadow $X_{O_r}=Y_{O_r}$; no common value vector is needed for equality terminalization.

Otherwise the proof does not expose the common value vector on $O_r$.  The only nonterminal information extracted from $O_r$ is finite role, scale, queue and incidence data.  Attempt to increment $M(i)$ for each $i\in O_r$.  If some increment would make $M(i)>Q$, then this $i$ is in the current actual split close-pair block, not merely in a possible support.  Therefore \Cref{lem:high-revisit} applies to that actual coordinate and gives alternative (ii).  The current block is charged to that hard exit and is not inherited as a low-revisit block.

If no cap is crossed, all incidences of $O_r$ are registered by the increments of $M$.  A bounded short window may be retained only as a bounded queue whose coordinate support is contained in already incremented marks.  Thus every coordinate identity used later is encoded in the retained state, and \Cref{lem:mark-compression} pays its entropy over the whole old-support epoch.  If the active signature changes, the previous signature's equality pattern has no meaning for the new labels.  Therefore the queue is either registered, terminalized, or charged to the signature-changing hard exit before being cleared; the child for the new signature is the quotient conditional law on the retained label.  This proves the alternatives and the assertion that no high-symbol old value vector is retained.
\end{proof}


\begin{lemma}[Reservoir marginalization]\label{lem:reservoir-marg}
Suppose a transition temporarily creates auxiliary sample labels only to certify a public event or realize conditioned copies.  If all future selectors and all retained state variables are measurable with respect to a retained set of labels and the public retained state, then replacing the full conditional law by its marginal on the retained labels preserves the law of all future selectors and terminal shadows involving retained labels.  It cannot increase the entropy of any retained key.
\end{lemma}

\begin{proof}
This is the elementary identity of marginal laws for events measurable in the retained coordinates.  Erased labels cannot be used by future selectors by hypothesis.  Deterministic marginalization is a data-processing map and cannot increase Shannon or Renyi entropy of a retained key.
\end{proof}

\begin{proposition}[Bounded reservoir]\label{prop:reservoir}
The number of active sample labels retained in a nonterminal state is bounded by an absolute constant.  Terminal label multiplicity satisfies
\[
        R(\Theta)\le T(\Theta)+O(1).
\]
\end{proposition}

\begin{proof}
All nonterminal modes retain only the active triple, pair, or bounded reservoir needed by the next transition.  Hash copying may temporarily create three conditioned copies, but after the copy tax is paid the auxiliary labels are erased by \Cref{lem:reservoir-marg}.  Every nonempty terminal sample label pins at least one coordinate, except for the fixed reservoir labels already present.
\end{proof}

\begin{proposition}[Scale-birth certificate]\label{prop:scale-birth}
Every new nonterminal active block $J$ is born with rank and scale balance paid by one of the following sources: initial active rank, parent active-rank drop, parent active-scale balance, fresh promotion, or old mark increments.  A hash token is assigned before a child born from its returned residual block is entered.  Terminal shadows and DRC selector certificates do not themselves create recurrence potential.
\end{proposition}

\begin{proof}
If $J$ is a fixed-fraction child of a parent active rank $m$, choose constants so that the parent rank drop pays the selector cost and the $B_s|J|$ scale deposit while retaining $B_{\rk}|J|$ in the child.  If $J$ is born from fresh coordinates, the fresh drop pays $(B_{\rk}+B_s)|J|$.  If it is born from old coordinates, the mark increments that carried $J$ forward pay the same amount through the old debit.  A token from a hash return is attached to the returned residual block and must be assigned to the hard exit that creates $J$, so the child is token-free.
\end{proof}

\section{The retained local transition theorem}\label{sec:local-transition}

This section proves the local theorem used in the assembly argument.  The alternatives are finite and mode-by-mode, and each transition is represented by one of the retained-state outputs accounted for below.

\begin{theorem}[Retained local transition theorem]\label{thm:local-transition}
At every nonterminal retained state generated by the rules below, one of the following retained exits is available, and the child law, if any, is the quotient law determined by the retained output.
\begin{enumerate}[label=\textup{(\roman*)}]
\item Inactive states use the collision dichotomy: terminal atom, low-collision near-sunflower residual activation, or public DRC close-pair activation.
\item Close-pair states use \Cref{prop:close-pair}.
\item Product-KL states use \Cref{prop:product-kl}.
\item Hash-window states use \Cref{thm:fixed-window} and \Cref{conv:hash-token}.
\item Residual states use \Cref{lem:fresh-promotion}, \Cref{prop:dense-residual}, \Cref{prop:bounded-residual}, and \Cref{prop:sparse-old}.
\item Old-window states use \Cref{prop:old-window}.
\item Reservoir and conditioned-copy labels are reduced by \Cref{prop:reservoir}; new active scales are funded by \Cref{prop:scale-birth}.
\end{enumerate}
In particular, no future selector depends on an erased analysis pattern, no hash key contains an erased old-pattern atom, every random old coordinate set carried forward is encoded by $M$ or a bounded queue/active block, and every pending hash token is assigned before a child state is entered.
\end{theorem}

\begin{proof}
We check the modes.  In an inactive state, if a codeword or coordinate atom is heavy it terminalizes.  If average coordinate collision is small, \Cref{lem:collision-split} gives a near-sunflower residual activation.  If average collision is large, the events $\{X_i=Y_i\ne Z_i\}$ have positive average density after heavy values are excluded, and \Cref{lem:drc} gives a public close-pair block.

Close-pair and product-KL states are covered by \Cref{prop:close-pair} and \Cref{prop:product-kl}.  Their nonterminal hash outputs are actual sample-coordinate laws under paid prefixes, and their terminal outputs are value/equality shadows or charged pruning branches.

For a hash window, the key is the current retained window transcript and paid set-events only.  By construction, erased old patterns are not part of the key.  If fewer than the target number of live coordinates are accumulated, the bounded support is immediately classified as fresh or old.  If enough live coordinates are accumulated, \Cref{thm:fixed-window} returns a residual block and \Cref{conv:hash-token} attaches exactly one bounded token.

For residual states, first apply the fresh support rule.  Once the branch is old-supported, the dense, bounded, and sparse alternatives are exhaustive according to whether $\E|W_L|$ is at least $\beta|L|$ and whether $|L|$ is bounded.  The dense and bounded cases are \Cref{prop:dense-residual} and \Cref{prop:bounded-residual}.  The sparse case is \Cref{prop:sparse-old}.  In every old nonterminal continuation the actual coordinates carried forward are contained in a mark increment or bounded queue.

Old-window states are \Cref{prop:old-window}.  Signature changes clear bounded queues before the new signature is used.  Reservoir marginalization and scale-birth funding are \Cref{prop:reservoir} and \Cref{prop:scale-birth}.  The retained-state conditions are therefore satisfied: terminal data is retained as shadows, monotone data is retained in finite state, pruning is charged, and analysis-only variables are quotiented before future choices are made.
\end{proof}

\section{Potential and global accounting}

Define the potential
\[
        \Phi=\Phi_{\rk}+\Phi_{\fr}+\Phi_{\old}+\Phi_{\scal}.
\]
The precise constants are chosen in a hierarchy.  It is enough for the proof that the following inequalities hold for fixed constants:
\begin{align*}
        B_f-B_o &> C_{\fr}+C_{\rm hash}+B_{\rk}+B_s+10,\\
        \omega_{\old} &> C_{\old}+B_{\rk}+B_s+C_{\rm hash}+10,\\
        B_o &> Q\omega_{\old}+C_{\rm blk}+C_{\rm aux}+10,\\
        B_{\rk}(1-\theta_{\rk}) &> (B_s+C_{\rk}+C_{\rm hash}+10)\theta_{\rk},\\
        B_s(1-\theta_{\scal})m_0 &> C_{\rm hash}+C_{\scal}.
\end{align*}
Here $C_{\rm blk}$ is the finite per-old-coordinate retained-state constant in \Cref{lem:mark-compression}.

\begin{proposition}[Branchwise retained charge]\label{prop:branch-charge}
For every retained branch of the finite transition tree,
\[
        \sum_{v\ {\rm reached}} c_v
        \le C\bigl(T(\Theta_{\rm val})+T(\Theta_{\rm eq})+R(\Theta)+\Delta\Phi\bigr),
\]
where $c_v$ is the retained selector charge at $v$, $\Delta\Phi$ is spent potential, and $C$ depends only on the fixed hierarchy.
\end{proposition}

\begin{proof}
Terminal value and equality exits are charged to their shadow sizes.  Fresh exits use disjoint promoted sets and the fresh drop.  Old exits use mark increments; by \Cref{lem:mark-compression} and the cap $Q$, all retained old non-value data is paid by $B_o|U|-\omega_{\old}\sum_i M(i)$ together with the per-coordinate state constant.  Scale descents are geometric after a scale is born; births are paid by \Cref{prop:scale-birth}.  Hash windows contribute one bounded token and \Cref{conv:hash-token} assigns it injectively to the first hard exit of the returned block.  Reservoir labels are bounded by \Cref{prop:reservoir}.  Summing these disjoint accounts gives the inequality.
\end{proof}

\begin{proposition}[Finite truncation]\label{prop:finite-truncation}
After bounded leaves are imposed below a fixed scale $m_0$, the retained transition system has finite truncations whose discarded tails are controlled by the pruning losses already charged.
\end{proposition}

\begin{proof}
It is enough to show that, after removing explicitly charged pruning branches, every nonterminal path has only finitely many retained changes before it reaches a terminal or bounded leaf.  Order the possible zero-pruning continuations by the lexicographic resource vector
\[
        \mathcal R=(N_{\rm fresh},N_{\rm old},N_{\rm hash},N_{\rm scale},N_{\rm queue}),
\]
where $N_{\rm fresh}$ is the number of still unpromoted coordinates, $N_{\rm old}=\sum_{i\in U}(Q-M(i))$ is remaining old mark capacity, $N_{\rm hash}$ is the bounded number of live coordinates still allowed before a window either returns or becomes a bounded-support exit, $N_{\rm scale}$ is the sum of dyadic scale levels above $m_0$ in all active blocks, and $N_{\rm queue}$ is the bounded unpaid queue capacity before an old window must register or clear.  The exact integer normalization of scale is fixed by the dyadic scale grid.

A fresh promotion strictly decreases $N_{\rm fresh}$ after depositing the old reserve.  A low-revisit old continuation strictly decreases $N_{\rm old}$; if this is impossible, the high-revisit exit is taken.  A hash-window step either adds one of fewer than $2h$ live coordinates, closes the bounded support, or returns a fixed-window residual block with a token; the token must be assigned before any child state or new hash window is entered, so it cannot create a cycle.  A deletion-trapped continuation decreases dyadic scale by the fixed factor, and any scale increase is a new birth paid by a rank, fresh, old, or scale certificate, so the paid birth is counted before the child is entered.  A bounded old queue can be extended only up to its fixed threshold; then it is registered, terminalized, or charged to a hard exit.

All deterministic tie-breaks are functions of the retained law and retained labels.  If none of the displayed resources changes and no pruning/terminal event occurs, the next retained state would be identical to the present one and the deterministic transition would repeat the same choice; by the transition rules this is declared a bounded leaf.  Thus finite depth truncations lose only explicitly pruned mass, whose normalization losses have already been recorded.  Letting the truncation depth tend to infinity preserves the entropy inequality by monotone convergence for the accumulated nonnegative pruning loss.
\end{proof}

\section{Assembly theorem}

\begin{theorem}[Retained-state assembly]\label{thm:assembly}
Choose constants as above.  There is a constant $D_0<\infty$, independent of $k$, the alphabet sizes, and $|\C|$, such that every balanced two-collision code of length $k$ has size at most $D_0^k$.  Consequently every family of sets of size at most $k$ with no three-petal sunflower has size at most $(eD_0)^k$.
\end{theorem}

\begin{proof}
Run the retained transition tree from the uniform law on the code.  By \Cref{lem:retained-entropy} and \Cref{prop:branch-charge}, after finite truncation and letting the truncation error tend to zero,
\[
        \HH(Z_*)+\E L_{\pr}
        \le C\,\E\bigl[T(\Theta_{\rm val})+T(\Theta_{\rm eq})+R(\Theta)+\Delta\Phi\bigr].
\]
Therefore some terminal retained outcome has original mass at least
\[
        \exp\{-C(T(\Theta_{\rm val})+T(\Theta_{\rm eq})+R(\Theta))-C\Delta\Phi\}.
\]
If the terminal output contains a nonempty value shadow, apply \Cref{thm:value-shadow}; if it contains only an equality shadow, apply \Cref{thm:equality-shadow}.  The potential compatibility follows from the construction of the branch accounts: spent potential plus inherited cylinder-child potential is at most the parent potential, because all hash tokens, old queues, and scale births are assigned before terminal counting.  Strengthen the induction to
\[
        |\C(\State)|\le D^{\rank(\State)}\exp(A\Phi(\State)).
\]
Choose $A$ and $D$ large enough to absorb the fixed constants.  The root potential is $O(k)$, so the bound is $D_0^k$ after increasing $D_0$.  The set-family statement follows from \Cref{cor:code-reduction}.
\end{proof}



\appendix

\section{Retained kernels and branchwise charges}\label{app:kernel-format}

This appendix expands the retained local transition theorem into checkable kernel statements.  The purpose is to make each local transition have the same form: input data, analysis partition, retained output, entropy charge, and quotient measurability.  All constants below are fixed after the threshold hierarchy is fixed; none depends on the alphabet sizes, on $k$, or on $|\C|$.

\begin{definition}[Retained kernel]\label{def:retained-kernel}
A retained kernel at a state $\State$ consists of the following data.
\begin{enumerate}[label=\textup{(K\arabic*)}]
\item A public state event $E_\State$ and a conditional law $\nu_\State$ on the active sample labels.
\item An analysis partition $\mathcal A=\{A_\omega\}$ of $E_\State$.  The atoms of $\mathcal A$ may use temporary equality patterns, KL prefixes, DRC witnesses, or old-window signatures.
\item A retained label map $Y$ on $\mathcal A$.  On a terminal atom, $Y$ is terminal shadow data.  On a nonterminal atom, $Y$ is one of the retained objects listed in \Cref{def:retained-data}: fresh promotion, old mark increment, bounded queue, active block, scale certificate, hash-window state, or pruning loss.
\item A quotient child law
\[
        \nu_y=\Law(\text{active samples}\mid E_\State\cap\{Y=y\}).
\]
Future selectors are deterministic functions of $y$, $\nu_y$, and the fixed thresholds.
\item A charge $c(\omega)$ satisfying
\[
        \HH(Y\mid E_\State)+\E[\ell_{\pr}\mid E_\State]
        \le \E[c(\omega)\mid E_\State].
\]
\end{enumerate}
\end{definition}

\begin{lemma}[Kernel composition]\label{lem:kernel-composition}
Suppose every nonterminal node of a finite retained tree is processed by a retained kernel in the sense of \Cref{def:retained-kernel}.  Then
\[
        \HH(Z_*)+\E L_{\pr}\le \E\sum_{v\ {\rm reached}} c_v.
\]
Moreover, if each nonterminal child is defined by the quotient law in \textup{(K4)}, no later selector can depend on an erased analysis atom.
\end{lemma}

\begin{proof}
Reveal only the retained labels $Y_v$ along the realised path.  The terminal object $Z_*$ is a deterministic function of this sequence.  Hence the chain rule gives
\[
        \HH(Z_*)\le \sum_v \Prob[v\ {\rm reached}]\HH(Y_v\mid E_v).
\]
Add the pruning losses.  By the local charge inequalities and monotone convergence for the finite truncations, the displayed bound follows.  The measurability assertion is immediate from the definition of the quotient child law: the sigma-algebra available at a child is generated by the retained label and the conditional law on the union of all analysis atoms with that label, not by the particular atom inside the union.
\end{proof}

\begin{proposition}[Branchwise charge table]\label{prop:appendix-charge-table}
Each retained output appearing in the local kernels below is charged by the following table.
\begin{center}
\begin{longtable}{>{\raggedright\arraybackslash}p{0.22\textwidth}>{\raggedright\arraybackslash}p{0.28\textwidth}>{\raggedright\arraybackslash}p{0.34\textwidth}}
\toprule
Output & Retained data & Charge source \\
\midrule
Value shadow & value vector on block $J$ for a sample label & $C_{\rm val}|J|+C_{\rm lab}$, paid by $T(\Theta_{\rm val})+R(\Theta)$ \\
Equality shadow & public or state-indexed block $J$ and two sample labels & $C_{\rm eq}|J|+C_{\rm lab}$, paid by $T(\Theta_{\rm eq})+R(\Theta)$ plus $\HH(S_*)$ if $J$ is state-indexed \\
Pruning & retained branch of relative mass $p$ & $-\log p$, recorded as $L_{\pr}$ \\
Hash coordinate & actual coordinate of one named sample label under a stopped key & at most $C_{\rm hash}$ per window token; window key entropy at most $B_{\rm win}$ \\
Fresh promotion & public set $F\subseteq [k]\setminus U$ & $C_{\fr}|F|$, paid by fresh potential drop \\
Old mark increment & support $K\subseteq U$ encoded by increasing $M(i)$ on $K$ & $C_{\old}|K|$, paid by $\omega_{\old}\sum_{i\in K}1$ and old reserve \\
Bounded old queue & at most $q_{\old}$ blocks of one signature and scale & $C_{\rm queue}|K|+O(1)$, paid by old reserve and cleared before signature change \\
Scale descent & active block replaced by smaller dyadic scale & $C_{\scal}|J|$, paid by scale potential drop after birth certificate \\
Scale birth & new active block at scale $s$ & $C_{\rm birth}|J|$, paid by rank/fresh/old source before child entry \\
Bounded leaf & all active blocks below $m_0$ & absorbed into the fixed base constant $D_0$ \\
\bottomrule
\end{longtable}
\end{center}
Consequently, after the constants satisfy the hierarchy in \Cref{prop:branch-charge}, the sum of retained charges is bounded by
\[
        C\bigl(T(\Theta_{\rm val})+T(\Theta_{\rm eq})+R(\Theta)+\Delta\Phi\bigr).
\]
\end{proposition}

\begin{proof}
Every row is either terminal data, monotone state data, pruning, or bounded base-case data.  Terminal rows are paid directly by the shadow descent coefficients.  Fresh coordinates are promoted once.  Old coordinates are charged only when their capped marks increase or when a bounded queue that contains them is retained; since $M(i)\le Q$, the total old retained entropy is at most the initial old reserve plus $Q$ increments per coordinate.  Hash windows have a single stopped key and a single token, which is assigned before a descendant state is entered.  Scale exits are geometric after a paid birth.  Summing the rows gives the branch inequality.
\end{proof}

\section{Close-pair and product-KL kernels}\label{app:close-product}

A close-pair state carries a public coordinate block $J$ and three active samples $(X,Y,Z)$ under the retained conditional law, with
\[
        X_J=Y_J=A,\qquad Z_J=B,
\]
where $B_j\ne A_j$ on the current close-pair event unless a terminal equality/value exit has already occurred.  Let $P=\Law(A_J,B_J)$ under the retained law.  Let $Q$ be the reference law obtained by taking the retained $A$-marginal and the retained $B$-marginal independently, then conditioning on the public side constraints already paid at the state.  All references to $P$ and $Q$ below are conditional on the current retained state.

\begin{lemma}[Actual-coordinate KL chain rule]\label{lem:actual-kl-chain}
List $J=(j_1,\ldots,j_m)$ in public order and set
\[
        V=(A_{j_1},B_{j_1},A_{j_2},B_{j_2},\ldots,A_{j_m},B_{j_m}).
\]
Then
\[
        \KL(P\|Q)=\sum_{t=1}^{2m}\E_P\KL\bigl(P(V_t\mid V_{<t})\,\|\,Q(V_t\mid V_{<t})\bigr).
\]
If the cumulative sum first crosses a threshold $b$, the crossing variable $V_t$ is an actual coordinate value of one of the active sample labels.
\end{lemma}

\begin{proof}
The map $(A_J,B_J)\mapsto V$ is a bijection of finite sample spaces.  Apply the ordinary chain rule for relative entropy to the ordered coordinates $V_1,\ldots,V_{2m}$.  Each $V_t$ is either $A_j=X_j=Y_j$ or $B_j=Z_j$ for an actual coordinate $j$ and an actual active sample label.
\end{proof}

\begin{lemma}[One-coordinate heavy/light routing]\label{lem:heavy-light-routing}
Let $W$ be an actual coordinate value of an active sample label under a paid prefix event.  Fix $0<\beta<1$.  Exactly one of the following retained exits is available.
\begin{enumerate}[label=\textup{(\alph*)}]
\item Some atom $\{W=a\}$ has conditional mass at least $\beta$.  The first such atom in the public order is a value-shadow exit, and the discarded complement is charged as pruning.
\item All atoms have mass below $\beta$.  Then $\Coll(W)=\sum_a\Prob(W=a)^2\le \beta$, so $W$ is a low-collision actual coordinate and enters the labelled hash window of its sample label under the current paid key.
\end{enumerate}
The retained charge is $O(1)$ plus the paid prefix/key entropy.
\end{lemma}

\begin{proof}
If a heavy atom exists, the public order chooses the first atom above threshold and the branch terminalizes with that value pin; the normalized complement is pruning.  If no heavy atom exists, then $\sum_a p_a^2\le (\max_a p_a)\sum_a p_a<\beta$.  The coordinate is not exposed by value and is appended to the fixed hash window for the corresponding sample label.
\end{proof}

\begin{lemma}[Positive-KL retained scan]\label{lem:positive-kl-retained-scan}
If $\KL(P\|Q)>B_{\rm cp}$, then the actual-coordinate scan of \Cref{lem:actual-kl-chain} produces, before the accumulated prefix budget exceeds $B_{\rm cp}$, either a terminal value shadow, a pruning branch, or a low-collision actual coordinate for a labelled hash window.  No retained output is a pair-alphabet coordinate or a variable of the reference law $Q$.
\end{lemma}

\begin{proof}
Run the cumulative scan in public order.  Prefixes whose description would exceed the remaining stopped budget are terminal value-shadow prefixes and are not inserted into a hash key.  Otherwise the first crossing variable is an actual coordinate value by \Cref{lem:actual-kl-chain}; apply \Cref{lem:heavy-light-routing}.  The retained data is either terminal shadow data, pruning loss, or a hash coordinate for a named active sample label.
\end{proof}

\begin{lemma}[Small-KL transfer]\label{lem:bounded-kl-transfer}
Let $P,Q$ be probability laws on a finite space and suppose $\KL(P\|Q)\le \delta^2/2$.  If $G$ is an event with $Q(G)\ge q$ and $0<\delta<q$, then
\[
        P(G)\ge q-\delta.
\]
In particular, if the close-pair hierarchy chooses the complement threshold $B_{\rm cp}\le q^2/8$, then $Q(G)\ge q$ implies $P(G)\ge q/2$.
\end{lemma}

\begin{proof}
By Pinsker's inequality,
\[
        \|P-Q\|_{\rm TV}\le \sqrt{\KL(P\|Q)/2}\le \delta/2.
\]
With the convention $\|P-Q\|_{\rm TV}=\sup_E |P(E)-Q(E)|$, this gives $P(G)\ge Q(G)-\delta/2\ge q-\delta$ after replacing $\delta$ by $2\delta$ in the fixed hierarchy.  The stated numerical form is obtained by choosing the hierarchy constants with slack.
\end{proof}

\begin{proposition}[Close-pair retained complement]\label{prop:appendix-close-complement}
In a close-pair state, suppose the following tests fail: terminal value/equality shadow, one-coordinate heavy/light routing, positive-KL scan with budget $B_{\rm cp}$, high $A$-side collision DRC, and terminal codeword atom.  Then the state has a retained residual output: a public near-sunflower block with exact split exposure, or a bounded leaf.  The retained charge is $O(|J|+B_{\rm cp})$, and the child law is the quotient law determined by the retained residual block.
\end{proposition}

\begin{proof}
Since the positive-KL scan fails, $\KL(P\|Q)\le B_{\rm cp}$.  Since high $A$-side collision fails, the $A$-marginal under $Q$ has average coordinate collision below the fixed threshold $\alpha$ on the public block $J$.  By the elementary collision-split lemma applied to three independent $A$-samples under $Q$, the event that the split set in $J$ has size at most $6\alpha |J|$ has $Q$-probability at least a fixed constant.  Removing the bounded bad events corresponding to non-distinct codewords and already-terminal atoms still leaves an event $G$ of $Q$-probability bounded below by a fixed constant, unless a terminal codeword atom holds.

Apply \Cref{lem:bounded-kl-transfer} with $B_{\rm cp}$ chosen below the fixed Pinsker threshold.  We obtain a retained event of $P$-mass bounded below by an absolute constant on which the projected block has the required near-sunflower split pattern.  The exact split exposure is retained only as the split block and role data needed for residual mode; any finer analysis atom is quotiented.  If $|J|<m_0$ the branch is a bounded leaf.  Otherwise the residual block is public, its split exposure costs $O(|J|)$, and the quotient child law is conditional on the union of all analysis atoms yielding the same retained residual block.
\end{proof}

\begin{corollary}[Close-pair kernel]\label{cor:appendix-close-kernel}
Every close-pair state admits one of the retained exits listed in \Cref{prop:close-pair}.  The local charge is bounded by a fixed multiple of the terminal shadow size, hash-window key charge, retained residual block size, or pruning loss.
\end{corollary}

\begin{proof}
Apply the tests in the order stated in the main text.  Terminal and heavy/light exits are handled by \Cref{lem:heavy-light-routing}; product-KL by \Cref{lem:positive-kl-retained-scan}; the complement by \Cref{prop:appendix-close-complement}.  Each nonterminal output is retained data, and all other analysis information is quotiented.
\end{proof}

\section{Fixed labelled hash windows}\label{app:hash}

\begin{lemma}[Window key independence]\label{lem:window-key-independence}
At the moment a coordinate is appended to a labelled hash window, the hash key consists only of the current retained state, the stopped window prefix, and paid set-events whose cumulative entropy is at most $B_{\rm win}$.  It contains no erased old-pattern atom and no old-window analysis signature unless that signature has been encoded into retained state.
\end{lemma}

\begin{proof}
Coordinates enter a window only through \Cref{lem:heavy-light-routing} or \Cref{lem:positive-kl-retained-scan}.  Both use the current quotient law and retained labels.  If an old-supported residual analysis was used earlier, its five-state pattern was quotiented before the child state was entered; by the retained kernel rule, later selectors are functions of the quotient law.  Therefore such a pattern cannot appear in the key.  If a pattern affects a later coordinate, it must do so through an old mark, bounded queue, active block, or token, all of which are retained and paid separately.
\end{proof}

\begin{lemma}[Renyi copy tax for a stopped key]\label{lem:renyi-key}
Let $K$ be a stopped public key with $\HH(K)\le B_{\rm win}$, and let $X$ be a sample label under the conditional law given $K$.  Three independent root samples can be conditioned to realize the same key at cost at most $O(B_{\rm win})$ in the terminal mass estimate.  More precisely, the matching probability is at least $\exp(-2B_{\rm win})$ after changing the absolute constant.
\end{lemma}

\begin{proof}
For a finite key, the order-three Renyi entropy satisfies $H_3(K)\le H(K)\le B_{\rm win}$.  The probability that three independent copies of the key coincide is $\sum_k p_k^3=\exp(-2H_3(K))\ge \exp(-2B_{\rm win})$.  Conditioning on a matching key gives three copies of the conditional sample law, and the copy cost is the negative logarithm of this probability.
\end{proof}

\begin{proposition}[Hash-window retained kernel]\label{prop:appendix-hash-kernel}
During one labelled hash window, either a terminal value/prefix shadow occurs, a bounded support is routed to fresh/old/scale output, or after $2h$ live low-collision coordinates the fixed-window copying lemma returns a residual near-sunflower block.  The total retained charge of the window is $O(B_{\rm win}+h)$ plus one bounded token assigned to the first hard exit of the returned residual block.
\end{proposition}

\begin{proof}
By \Cref{lem:window-key-independence}, the key has entropy at most $B_{\rm win}$ and contains no erased analysis atom.  A coordinate can be appended only while its prefix, tail/pruning, and production transcript fit into the stopped budget.  If the next addition would exceed the budget, the existing prefix terminalizes as a value shadow or the current bounded support exits to fresh/old/scale processing.  Nonterminal live coordinates have no atom of mass at least $\beta$, by the survival scan, so their coordinate collisions are at most $\beta$.

If $2h$ distinct live coordinates are present, choose any public $h$ surviving coordinates after the deterministic tie-break.  By \Cref{lem:renyi-key}, three conditioned copies with the same key cost $O(B_{\rm win})$.  On the $h$ low-collision coordinates the probability of a bad split pattern is $O(\beta h)$; choosing $\beta h$ small gives positive mass for a near-sunflower block.  Exposing the exact split set costs $O(h)$.  The returned block carries exactly one bounded token, which is assigned before a descendant state or new window is entered.
\end{proof}

\section{Residual dense, bounded, and sparse kernels}\label{app:residual}

A residual state carries a public near-sunflower block $L$ and exact split exposure for an active triple.  Let $U$ be the old support and let $W_L(X,Y,Z)$ be the split witness inside $L$.

\begin{lemma}[Fresh support kernel]\label{lem:appendix-fresh-kernel}
Let
\[
        F=\{i\in L\setminus U:\Prob[i\in W_L]>0\}
\]
under the current quotient law.  If $|F|>q_{\fr}$, promote $F$ to $U$ and charge the fresh potential.  If $|F|\le q_{\fr}$, expose the five-state pattern on $F$ as an analysis partition; each nonterminal child either promotes a nonempty fresh split part or has all residual split witnesses contained in $U$.  The only retained data is the promoted set, and the cost is $O(|F|)$.
\end{lemma}

\begin{proof}
The set $F$ is determined by the quotient law.  If it is large, promotion is public and each coordinate is promoted once.  If it is bounded, the five-state pattern on $F$ has $5^{|F|}$ possibilities, hence entropy $O(|F|)$.  On an atom with a fresh split coordinate, retain the corresponding promoted split part.  On atoms with no fresh split coordinate, quotient the pattern and continue with old-supported residual witnesses.  Future choices are functions of the retained promoted set or the old-supported quotient law.
\end{proof}

\begin{lemma}[Dense old-supported kernel]\label{lem:appendix-dense-old}
Assume $W_L\subseteq U$ and $\E|W_L|\ge \beta |L|$.  Then there is a public old block $K\subseteq U$ of fixed size $s=s(\beta)$ such that
\[
        \Prob[K\subseteq W_L]\ge \exp(-O(s)).
\]
On this branch, $K$ is either terminalized as an equality shadow, routed to an old-window block, or encoded as old mark increments.  The retained charge is $O(s)$.
\end{lemma}

\begin{proof}
For $i\in L$, let $E_i=\{i\in W_L\}$.  The average of $\Prob(E_i)$ over $L$ is at least $\beta$.  Apply the DRC lemma to the events $E_i$ with target size $s$; the first valid block in the public order is $K$.  Conditioning on $K\subseteq W_L$ gives an actual old split block.  A majority of the split roles gives a close-pair subblock of size at least $|K|/3$.  If the close-pair event has terminal mass, use equality shadow.  Otherwise retain the block only by increasing $M$ on its support or by placing it in the bounded old queue.  All other analysis data is quotiented.
\end{proof}

\begin{lemma}[Bounded residual kernel]\label{lem:appendix-bounded-residual}
If the residual support $L$ has bounded size, then exposing the complete five-state pattern on $L$ costs $O(1)$ and the branch is routed to a terminal shadow, bounded leaf, fresh promotion, or old mark increment.
\end{lemma}

\begin{proof}
There are $5^{|L|}=O(1)$ pattern atoms.  Each atom determines the split set and role.  Terminal atoms are retained as shadows; nonterminal old blocks are encoded through $M$ or the bounded queue; fresh blocks are promoted.  Since $|L|$ is bounded, the entropy is absorbed into the finite local charge.
\end{proof}

\begin{lemma}[Sparse old-supported retained kernel]\label{lem:appendix-sparse-old}
Assume $W_L\subseteq U$ and $\E|W_L|<\beta|L|$.  Let
\[
        O=\{u\in U:\Prob[u\in W_L]>0\}.
\]
Use the five-state equality pattern on $O$ only as an analysis partition.  On each atom, the actual split set $K=W_L$ is determined.  Then exactly one of the following retained exits is taken:
\begin{enumerate}[label=\textup{(\alph*)}]
\item $K=\emptyset$, in which case the branch is pruned or declared a bounded residual leaf according to the retained mass;
\item a public or state-indexed close-pair subblock $K_0\subseteq K$ terminalizes as an equality shadow;
\item retaining $K_0$ would cross the cap $Q$ at some coordinate, so the high-revisit kernel of \Cref{lem:appendix-high-revisit} is applied;
\item otherwise $K_0$ is encoded by increasing $M(i)$ on $i\in K_0$, possibly together with bounded role/queue marks.
\end{enumerate}
The child law is the quotient law over all analysis atoms giving the same retained output.  The retained old entropy over the whole branch is $O(|U|)$.
\end{lemma}

\begin{proof}
The set $O$ is public relative to the quotient law.  The five-state pattern on $O$ may have large entropy, so it is not a retained label.  It is used only to determine, on each atom, the actual split set $K$ and the majority close-pair role.  If a terminal equality shadow is taken, the block identity is either public or included in the terminal retained state; the state-indexed shadow descent pays its identity through $\HH(S_*)$.  If the branch is nonterminal, every coordinate of the block used later is marked by an increment of $M$ or placed in a bounded queue.  Therefore the retained label records only a subset through capped marks and finite role data.  By \Cref{lem:mark-compression}, the number of possible retained old states is at most $\exp(O(|U|))$.  All other pattern entries are quotiented, and future choices are computed from the conditional law on the union of atoms with the same retained label.
\end{proof}

\section{Old-window and high-revisit kernels}\label{app:old-window}

\begin{lemma}[High-revisit one-coordinate routing]\label{lem:appendix-high-revisit}
Suppose an old coordinate $i\in U$ would be marked more than $Q$ times by actual close-pair blocks.  Let $A_i$ be the common value on the current actual pair at $i$, under the current retained law.  Then one of the following exits is available:
\begin{enumerate}[label=\textup{(\alph*)}]
\item a heavy atom of $A_i$ has mass at least $\beta$, giving a terminal value shadow and charged pruning of the complement;
\item all atoms have mass below $\beta$, giving $\Coll(A_i)\le \beta$, so $i$ is a low-collision actual coordinate for the labelled hash window;
\item the remaining branch has bounded retained mass and is charged as pruning or bounded leaf.
\end{enumerate}
No high-symbol value vector is retained in a nonterminal child.
\end{lemma}

\begin{proof}
Apply the heavy/light routing lemma to the one-coordinate law of $A_i$.  A heavy atom terminalizes; on the light complement the collision is at most $\beta$ and the coordinate is appended to the appropriate labelled hash window under the current retained key.  If the conditioning event needed to isolate the current close-pair block has too small retained mass, the branch is charged as pruning.  The value of $A_i$ is never retained unless it is terminal.
\end{proof}

\begin{proposition}[Old-window retained kernel]\label{prop:appendix-old-window}
An old-window state with one sample-label signature and one dyadic scale has a retained exit of one of the following forms:
\begin{enumerate}[label=\textup{(\roman*)}]
\item terminal value or equality shadow;
\item high-revisit routing by \Cref{lem:appendix-high-revisit};
\item low-revisit continuation, encoded by old mark increments and bounded queue data;
\item signature change, after the current queue is registered, terminalized, pruned, or cleared by a hard exit.
\end{enumerate}
The retained child contains no old value vector and no discarded signature-specific equality pattern.
\end{proposition}

\begin{proof}
Run the old blocks in public order.  If a block has terminal equality mass, stop by equality shadow; if a low-entropy value vector is selected, stop by value shadow.  Otherwise each actual old incidence is tested against the cap $Q$.  Crossing the cap invokes \Cref{lem:appendix-high-revisit}.  If no cap is crossed, increment $M$ on the actual old coordinates used later and store only bounded role/queue information.  Before a new sample-label signature is used, the current queue must be empty or have taken one of the hard exits above.  Thus no child state contains two unmerged signature patterns and no future selector can use an erased pattern atom.
\end{proof}

\section{Reservoir, scale birth, and finite termination}\label{app:termination}

\begin{lemma}[Reservoir reduction]\label{lem:appendix-reservoir}
Auxiliary conditioned-copy labels and temporary DRC samples are not retained beyond the local transition in which they are used.  Any sample label that remains active in a nonterminal child belongs to a reservoir of fixed size $r_*$, and the retained label identifying it has entropy $O_{r_*}(1)$ per transition.
\end{lemma}

\begin{proof}
Conditioned copies are realised inside the local transition and then either terminalize, produce a residual block, or are discarded.  The retained child law is the marginal law of the bounded set of active labels that remain.  Since the reservoir size is fixed, relabelling and tie-breaking among active labels costs only a constant depending on $r_*$.  Discarded labels are not available to future selectors.
\end{proof}

\begin{lemma}[Scale birth certificate]\label{lem:appendix-scale-birth}
Every new nonterminal active block is born with a certificate paid before the child state is entered.  The certificate is one of: active-rank drop, fresh promotion, old mark increment, scale descent, or hash-token assignment.  Hence scale potential is never increased for free.
\end{lemma}

\begin{proof}
A DRC child of a rank-$m$ active block has fixed fractional size, so the parent rank drop pays the selector and deposits the child scale balance.  A child supported on newly promoted coordinates is paid by fresh potential.  A child supported on old coordinates is paid by the mark increments or bounded queue that carries its coordinate identity.  A hash-returned residual child carries one token, and the transition rules assign that token to the first hard exit before any descendant state is entered.  These cases exhaust active block births.
\end{proof}

\begin{proposition}[No zero-cost retained loop]\label{prop:appendix-no-loop}
After bounded leaves below $m_0$ are imposed, there is no infinite nonterminal path on which no terminal shadow, pruning charge, fresh promotion, old mark increment, hash-window token assignment, scale descent, or paid scale birth occurs.
\end{proposition}

\begin{proof}
Use the lexicographic resource vector
\[
        (N_{\fr},N_{\old},N_{\mathrm{hash}},N_{\mathrm{scale}},N_{\mathrm{queue}}),
\]
where $N_{\fr}$ counts unpromoted fresh coordinates, $N_{\old}=\sum_{i\in U}(Q-M(i))$, $N_{\mathrm{hash}}$ is the remaining capacity in the current stopped hash window before return or bounded-support exit, $N_{\mathrm{scale}}$ is the total dyadic scale level above $m_0$, and $N_{\mathrm{queue}}$ is remaining bounded old-queue capacity.  Every nonterminal retained continuation strictly decreases one entry after all earlier entries are fixed, unless a paid birth resets a later scale entry; such a birth is itself charged before the child state is entered.  If none of the entries changes and no terminal/pruning event occurs, the deterministic transition from the same retained law and same retained labels repeats identically.  The rules then declare the state a bounded leaf.  Therefore finite truncations lose only explicitly pruned mass.
\end{proof}

\section{Appendix conclusion}\label{app:conclusion}

Combining \Cref{lem:kernel-composition}, the charge table of \Cref{prop:appendix-charge-table}, the close-pair and product-KL kernels of \Cref{app:close-product}, the hash kernel of \Cref{app:hash}, the residual kernels of \Cref{app:residual}, the old-window kernel of \Cref{app:old-window}, and the termination lemmas of \Cref{app:termination} gives the retained local transition theorem used in \Cref{thm:local-transition}.  The appendices are written so that each temporary analysis object is accounted for in exactly one of four ways: terminal shadow, retained finite state, pruning loss, or quotient-erased analysis data.

\begin{thebibliography}{9}
\bibitem{ALWZ}
R. Alweiss, S. Lovett, K. Wu, and J. Zhang, \emph{Improved bounds for the sunflower lemma}, Annals of Mathematics 194 (2021), 795--815.

\bibitem{ER}
P. Erd\H{o}s and R. Rado, \emph{Intersection theorems for systems of sets}, J. London Math. Soc. 35 (1960), 85--90.

\bibitem{Rao}
A. Rao, \emph{Coding for sunflowers}, Discrete Analysis 2020:2 (2020), 8 pp.

\bibitem{Tao}
T. Tao, \emph{The sunflower lemma via Shannon entropy}, blog note, 2020.
\end{thebibliography}

\end{document}
